% ECONOMICS_PROBLEM: {"id": "H41-557", "title": "Identity, ethnicity and redistribution/public goods", "jel_code": "H41", "source_id": "OP-0129"}
% !TeX program = xelatex
\documentclass[11pt,a4paper]{article}
\usepackage[margin=23mm]{geometry}
\usepackage{fontspec}
\setmainfont{Latin Modern Roman}
\usepackage{amsmath,amssymb,mathtools}
\usepackage{xcolor,enumitem,booktabs,longtable,array,xurl,hyperref}
\definecolor{muted}{HTML}{53636C}
\hypersetup{colorlinks=true,urlcolor=blue,linkcolor=blue}
\setlength{\parindent}{0pt}
\setlength{\parskip}{5pt}
\newcommand{\R}{\mathbb R}
\newcommand{\E}{\mathbb E}
\newcommand{\N}{\mathbb N}
\newcommand{\ind}{\mathbf 1}
\newcommand{\Var}{\operatorname{Var}}
\newcommand{\Cov}{\operatorname{Cov}}
\newcommand{\supp}{\operatorname{supp}}
\DeclareMathOperator*{\argmin}{arg\,min}
\DeclareMathOperator*{\argmax}{arg\,max}
\newcommand{\entrymeta}[3]{{\sffamily\small\color{muted}\textbf{\texttt{#1}}\quad|\quad #2\par #3\par}}
\newcommand{\Problemref}[1]{\texttt{#1}}
\newcommand{\JELref}[2]{\texttt{#2} (JEL \texttt{#1})}
\begin{document}
\section*{Identity, ethnicity and redistribution/public goods}
\textbf{Problem ID:} \texttt{H41-557}\quad \textbf{JEL:} \texttt{H41}\par
% BEGIN ECONOMICS STATEMENT
\label{problem:OP-0129}
\label{atlas:prob:Q9}
{\sffamily\small\color{muted}\textbf{Q9} | Political economy | Editorial research problem family\par}
\subsection*{Definitions and assumptions}
Fix a population $N$ with immutable recorded background characteristics $G_i$, and a feasible institutional/salience policy $a$ that does not purport to randomize ancestry. Individual identity $I_i(a)$, public-good contribution $c_i(a)\geq0$, tax support $T_i(a)$ and behavior are endogenous. A public-good technology $q(a)=F(\sum_i c_i(a),a)$ and collection costs are specified. Utilities $u_i(x_i,q,I_i,G_i,a)$ define a game; the model class includes identity formation, social monitoring, incomes and an equilibrium selection rule. Set fiscal feasibility and a declared welfare function $W(a)$ with explicit interpersonal weights.
\subsection*{Formal research question}
Identify or bound population mean contribution and tax-support effects of $a$ versus $a_0$, and distinguish mechanisms through feasible interventions on communication, enforcement, group salience or collective institutions:
\[
\theta_c(a,a_0)=\frac1{|N|}\sum_{i\in N}\mathbb E[c_i(a)-c_i(a_0)].
\]
Determine which designs or model restrictions separate preference changes, production/coordination technology and equilibrium selection. When ranking institutions, characterize the feasible maximizers or near-maximizers of $W(a)$, retaining distributional and identity effects.
\subsection*{Interpretation and scope}
Cross-sectional ethnic fractionalization is a population descriptor rather than a manipulable treatment. A causal comparison must identify an institution, boundary assignment, salience change or composition intervention with well-defined versions. Contributions in a laboratory game need not equal sustained tax compliance or public-service outcomes. Endogenous identity also prevents an unqualified interpretation of identity-conditioned comparisons as causal effects of an exogenous personal attribute.
\subsection*{Source audit and status}
All three original citations are verified. Alesina--Baqir--Easterly supplies municipal evidence, Habyarimana et al. compares mechanisms experimentally, and Shayo endogenizes identity in equilibrium. These sources already provide partial explanations. The remaining institutional and transportability target is editorial, with no universal claim that diversity reduces provision in every setting.
\subsection*{References}
{\footnotesize\raggedright
\begin{enumerate}[label={\textbf{[S\arabic*]}},leftmargin=3em]
\item Alberto Alesina, Reza Baqir, and William Easterly (1999). \emph{Public Goods and Ethnic Divisions}. Quarterly Journal of Economics 114(4), 1243--1284.\par
\textit{Locator and source role:} Abstract and municipal public-goods evidence; does not make ethnicity an exogenous manipulable treatment.\par
\url{https://academic.oup.com/qje/article-abstract/114/4/1243/1934019}
\item James Habyarimana, Macartan Humphreys, Daniel N. Posner, and Jeremy M. Weinstein (2007). \emph{Why Does Ethnic Diversity Undermine Public Goods Provision?}. American Political Science Review 101(4), 709--725.\par
\textit{Locator and source role:} Abstract and experimental mechanism comparisons in Kampala; distinguishes preferences, technology and strategy selection.\par
\url{https://www.cambridge.org/core/journals/american-political-science-review/article/abs/why-does-ethnic-diversity-undermine-public-goods-provision/36C285D07DC0DE604BFC123F53060DF6}
\item Moses Shayo (2009). \emph{A Model of Social Identity with an Application to Political Economy: Nation, Class, and Redistribution}. American Political Science Review 103(2), 147--174.\par
\textit{Locator and source role:} Abstract and endogenous identity equilibrium; contextual rather than universal causal effects of ethnic composition.\par
\url{https://doi.org/10.1017/S0003055409090194}
\end{enumerate}
}

\subsection*{Common conventions: Source mathematical conventions}

Unless an entry states otherwise, agent and firm sets are finite. State and action spaces are measurable, strategies and outcome maps are measurable, probability laws are specified, and expectations exist for the targets under discussion. Infinite-horizon discount factors lie in $(0,1)$ and discounted objectives are finite. Norms, tolerances and complexity input sizes must be specified for computational comparisons. Constants or primitives that appear locally are inputs, not empirical facts asserted by this catalogue.

\textbf{Identification.} A statistical model $m\in\mathcal M$ implies an observable law $P_m$ and target $\theta(m)$. At law $P$, its identified set is
\[
\Theta(P;\mathcal M)=\{\theta(m):m\in\mathcal M,\ P_m=P\}.
\]
Point identification holds when this set is a singleton, or equivalently when $P_{m_1}=P_{m_2}$ implies $\theta(m_1)=\theta(m_2)$ for all compatible models. Sharp bounds mean bounds attained, or approached when the boundary is not attained, by compatible models. Writing this set is a definition, not a solution: the substantive task is to establish informative restrictions and derive or estimate the set. An unrestricted-confounding model may give only trivial bounds.

\textbf{Inference.} Identification, estimability and valid uncertainty quantification are separate questions. Fix $\alpha\in(0,1)$. A confidence procedure $C_n$ over a declared law class $\mathcal P$ has asymptotic uniform coverage at level $1-\alpha$ when
\[
\liminf_{n\to\infty}\inf_{P\in\mathcal P}\Pr_P\{\theta(P)\in C_n\}\geq1-\alpha.
\]
For set-identified targets the coverage event must be defined separately, for example coverage of the full identified set. If an entry uses identification-indexed models instead of uniquely indexed laws, the same coverage convention is applied over those models. Pointwise limits do not establish uniform validity, and asymptotic validity does not guarantee finite-sample precision. Sample sizes, dependence structures and weak-identification sequences are part of each application.

\textbf{Counterfactuals.} $Y_i(a)$ is a potential outcome under a specified intervention, including its timing, implementation and versions. It is not $Y_i$ conditional on voluntarily choosing $a$. A full intervention vector permits interference; shorthand scalar treatment notation does not silently impose no interference. Assignment, selection, attrition, anticipation and measurement must be specified. A claim about an instrument's local effect is not automatically a population policy effect.

\textbf{Economic equilibrium.} A counterfactual model includes best responses, constraints, feasibility, market clearing, state transitions and expectations consistent with the stated information. When several equilibria exist, either a declared selector is an input or the target is set-valued. Comparisons across policy regimes must use the stated selection convention. Reduced-form relationships are not presumed invariant after an equilibrium-changing intervention.

\textbf{Optimization and welfare.} Optimization is over the stated feasible set. If attainment is not established, $\sup$ or $\inf$ and $\varepsilon$-optimal choices replace an asserted maximizer or minimizer. For example, with value $V=\sup_{a\in\mathcal A}W(a)<\infty$, an $\varepsilon$-optimal set is $\{a\in\mathcal A:W(a)\geq V-\varepsilon\}$ for $\varepsilon>0$. Benefits, costs, risk and health or environmental outcomes can be aggregated only using declared units and conversion weights. Interpersonal utility weights, the treatment of fiscal transfers and foreign profits, discounting, and the behavioral welfare criterion are explicit inputs. A comparative-static derivative additionally requires existence and appropriate differentiability of the selected counterfactual.

\textbf{Sources.} Local labels [S1], [S2], and so forth restart in every question. Locators identify the relevant abstract, model, section or result at the level actually checked; a bibliographic or abstract-level verification is not represented as inspection of an entire proof. Working papers are distinguished from later publications and changing versions. Source summaries are paraphrased. All URL checks and editorial formulations refer to the audit date stated above.
% END ECONOMICS STATEMENT
\end{document}
